\section{Weighted free pro-2 groups}
\label{sec:weighted}

We will construct an infinite quotient by comparing the weighted number of generators with the weighted number of relations. This section records the precise version of the Golod--Shafarevich argument that we use, including convergence for infinitely many relators.

\subsection{Magnus degrees and automorphisms}

A pro-2 group is an inverse limit of finite groups of order a power of two. The free pro-2 group on a finite set is the pro-2 completion of the discrete free group on that set. Let $F_d$ be free pro-2 on $g_1,\ldots,g_d$, and assign positive integer weights $w_1,\ldots,w_d$. Its Magnus expansion sends
\[
 g_i\longmapsto1+U_i
 \quad\text{in}\quad
 A=\F_2\langle\!\langle U_1,\ldots,U_d\rangle\!\rangle,
 \qquad \deg U_i=w_i.
\]
Here $A$ is the algebra of formal noncommutative power series, with the filtration $A_{\ge n}$ by weighted degree at least $n$. For $g\ne1$, let $\nu(g)$ be the least degree in $g-1$, and put $\nu(1)=\infty$.

For clarity, this expansion is faithful and gives the pro-2 topology. Each truncated algebra is finite, and its group $1+A_{>0}/A_{\ge n}$ is a finite 2-group. Conversely the augmentation ideal of $\F_2[P]$ is nilpotent for every finite 2-group $P$. To see this, induct on $|P|$, choosing a central involution $z$ when $P\ne1$. The kernel of $\F_2[P]\to\F_2[P/\langle z\rangle]$ is generated by $z-1$ and has square zero. Nilpotence follows by induction. Thus every finite 2-quotient of the free group factors through some Magnus truncation. This proves faithfulness on the pro-2 completion and cofinality of the truncations. Positive finite weights give the same topology as ordinary degree.

Multiplication of series gives
\begin{equation}\label{eq:degree-rules}
 \nu(gh)\ge\min\{\nu(g),\nu(h)\},\qquad
 \nu([g,h])\ge\nu(g)+\nu(h),\qquad
 \nu(g^{2^e})\ge2^e\nu(g).
\end{equation}
We use either commutator convention, since only its degree matters. The last inequality follows from $g^{2^e}-1=(g-1)^{2^e}$ in characteristic two.

\begin{lemma}[Substitution]\label{lem:substitution}
Suppose an automorphism $\alpha$ of $F_d$ satisfies
\[
 \nu\bigl(\alpha(g_i)g_i^{-1}\bigr)\ge w_i+a
 \quad(1\le i\le d)
\]
for an integer $a\ge1$. Then its continuous action on $A$ satisfies
\[
 (\alpha-\id)(A_{\ge n})\subset A_{\ge n+a}.
\]
More generally, substitution of degree-raising operators for Magnus variables is well defined whenever the operator assigned to a variable of weight $w$ raises degree by at least $w$.
\end{lemma}

\begin{proof}
The hypothesis says that $\alpha(U_i)-U_i$ has degree at least $w_i+a$. Expand the difference between a monomial and its substituted image. Every summand replaces at least one factor by an error of at least $a$ additional degrees. The estimate follows for polynomials and then for series by completeness. For operator substitution, a word of weighted degree $n$ raises degree by at least $n$; hence only finitely many words contribute in each finite truncation. Multiplication and composition respect these substitutions.
\end{proof}

We also use the following basis criterion. The \emph{Frattini quotient} of a finitely generated pro-2 group $G$ is
\[
 G/\overline{G^2[G,G]},
\]
an $\F_2$-vector space. Elements whose images span this space topologically generate $G$: otherwise their closed generated subgroup has proper image in a finite 2-quotient, and is contained in a maximal subgroup of index two there, contradicting the span. Consequently, $d$ elements of $F_d$ whose Frattini images form a basis are again a free pro-2 basis. Indeed, they define a surjective endomorphism of $F_d$, and every finitely generated profinite group is Hopfian. For the latter assertion, a surjective endomorphism permutes, by inverse image, the finite set of open normal subgroups of any fixed index. Its kernel therefore lies in every open normal subgroup and is trivial.

\subsection{The infinitude inequality}

The infinitude method originates in the work of \citet{GolodShafarevich1964}; \citet{Vinberg1965} developed its filtered form. For the weighted completed-algebra inequality and its pro-$p$ consequence, see \citet[Theorem~2.1 and Corollary~2.2]{Ershov2011}. The more general language of positive weighted deficiency is developed by \citet[arXiv version, Corollary~4.4]{ErshovJaikinZapirain2013}. We include the filtered proof needed here to specify exactly the completion and the infinite-relator convention.

\begin{lemma}[Weighted infinitude criterion]\label{lem:weighted-gs}
Let $F_d$ have the weighted basis above, and let $(r_\lambda)$ be a finite or countable family of relators with positive integer bounds $v_\lambda$ satisfying
\[
 \nu(r_\lambda)\ge v_\lambda\ge1.
\]
Assume only finitely many $v_\lambda$ lie below any fixed bound. If, for some $0<t<1$, the relator series converges and
\begin{equation}\label{eq:gs-criterion}
 1-\sum_{i=1}^d t^{w_i}+\sum_\lambda t^{v_\lambda}<0,
\end{equation}
then the quotient of $F_d$ by the closed normal subgroup generated by all $r_\lambda$ is infinite.
\end{lemma}

\begin{proof}
Let $I\subset A$ be the closed two-sided ideal generated by $f_\lambda=r_\lambda-1$, and let $R=A/I$ with its quotient filtration. Put
\[
 C(n)=\dim_{\F_2}R/R_{\ge n}\quad(n\ge1),
 \qquad C(n)=0\quad(n\le0).
\]
First-letter decomposition gives a surjective map
\[
 \bigoplus_i R/R_{\ge n-w_i}\longrightarrow
 R_{>0}/R_{\ge n},\qquad
 (b_i)_i\longmapsto\sum_i U_i b_i.
\]
Its kernel has dimension at most $\sum_\lambda C(n-v_\lambda)$. Here is the filtered justification. Write $f_\lambda=\sum_i U_i f_{i\lambda}$. A vector in the kernel can be lifted to series whose first-letter sum lies in $I+A_{\ge n}$. Subtracting the first-letter decomposition of the degree-at-least-$n$ term does not change the truncated vector, so the sum may be taken in $I$. For a product $a f_\lambda h$, the positive-degree part of $a$ contributes first-letter coefficients that vanish in $R$; its constant part contributes
\[
 (\overline{f_{i\lambda}})_i\,\bar h.
\]
This vector has shifted degree at least $v_\lambda$. Only $\bar h$ modulo $R_{\ge n-v_\lambda}$ matters. These vectors therefore span a space of dimension at most $\sum_\lambda C(n-v_\lambda)$. Closure adds nothing to a span in a finite-dimensional truncation. Since $\dim R_{>0}/R_{\ge n}=C(n)-1$, we obtain
\begin{equation}\label{eq:cumulative-gs}
 C(n)-\sum_i C(n-w_i)+\sum_\lambda C(n-v_\lambda)\ge1.
\end{equation}

Suppose the stated pro-2 quotient were finite, of order $m$. Its image spans every truncated algebra $R/R_{\ge n}$: the variables are the images of $g_i-1$, and their products are linear combinations of group elements. The group map to each truncated unit group kills the relators, hence factors through that quotient. Thus $C(n)\le m$ for all $n$.
Multiply \eqref{eq:cumulative-gs} by $t^{n-1}$ and sum over $n\ge1$. Boundedness of $C(n)$ and convergence of the relator series justify the sums. This gives
\[
 \left(1-\sum_i t^{w_i}+\sum_\lambda t^{v_\lambda}\right)
 \sum_{n\ge1}C(n)t^{n-1}\ge\frac1{1-t},
\]
contradicting \eqref{eq:gs-criterion}.
\end{proof}

The use of lower bounds $v_\lambda$ makes the criterion convenient: increasing an actual relator degree can only improve the inequality. In the construction below, the inexpensive relations impose a geometric involution. The remaining relations will be pushed arbitrarily far into the filtration.
